\documentclass{situationist-preprint}

\title{Trying, Analyzed}
\msid{ACT-11-b}
\msfamily{ACT-11}
\msdiscipline{Action, Agency and Free Will}
\mstype{Companion manuscript}
\mskeywords{trying; intention; volition; aim; guidance; intending to try}

\begin{document}
\maketitle

\begin{abstract}
ACT-11-a argues that one cannot intend what one believes one will not do, but can intend to try. This requires an account of trying on which trying does not itself involve intending the outcome. We give one. To try to $A$ is to perform some act, bodily or mental, that one believes to be a way of $A$-ing or of raising the chance that one $A$s, under the guidance of $A$-ing as its standard of success. Guidance is weaker than intention, and the account shows why: if trying required intending the outcome, Theorem~3.2 of ACT-11-a would contradict its Theorem~3.1.
\end{abstract}

\section{The analysis}

Definition~2.2 of ACT-11-a, following \citet{OShaughnessy1973} and \citet{Hornsby1980}, takes trying to $A$ to be doing what one believes is a way of $A$-ing, or of making it more likely that one $A$s, with the aim of $A$-ing. We make each element explicit.

\begin{definition}[Trying, analyzed]\label{d:try}
An agent \emph{tries} to $A$ at $t$ if and only if (i) she performs at $t$ some act $B$, which may be a bodily movement or a mental act such as an act of will; (ii) she believes that $B$ is, or may be, a way of $A$-ing or of raising the chance that she $A$s; and (iii) her $B$-ing is guided by $A$-ing, in that she monitors its progress toward $A$-ing and is disposed to adjust or repeat it so as to $A$.
\end{definition}

Clause~(i) allows the paralyzed patient who, not yet knowing of her paralysis, tries to raise her arm: she performs an act of will, believed to be a way of raising it, and fails \citep{OShaughnessy1973}. Clause~(ii) excludes the agent who does something she knows to be irrelevant to $A$-ing: one who whistles while wishing the bus would come is not trying to make it come. Clause~(iii) distinguishes trying from going through the motions: the swimmer who swims at a pace she knows will not win, indifferent to the outcome, is not trying to win, though she does what is in some sense a way of winning.

\begin{proposition}\label{p:guide}
The guidance of clause~(iii) cannot be an intention to $A$.
\end{proposition}

\begin{proof}
By Theorem~3.2 of ACT-11-a, an agent may intend to try to $A$, and so may try, while believing that she will not $A$. If trying to $A$ required intending to $A$, such an agent would intend to $A$ while believing she will not, which Theorem~3.1 of ACT-11-a excludes. So the aim in trying is not an intention to $A$.
\end{proof}

Guidance by $A$-ing is a matter of $A$-ing's functioning as the standard by which the agent assesses and adjusts her conduct. It requires that she want to $A$; it does not require that she settle the question whether she will $A$, which, by Principle~P1 of ACT-11-a, she cannot do while believing she will not. What she settles is a different question, whether to try, and that is the content of the intention that Theorem~3.2 allows.

\begin{proposition}\label{p:ubiq}
Every intentional $A$-ing involves trying to $A$, but not every trying to $A$ involves intending to $A$.
\end{proposition}

\begin{proof}
An agent who intentionally $A$s does something she believes to be a way of $A$-ing, under the guidance of $A$-ing; so she satisfies clauses~(i)--(iii). The converse fails by Proposition~\ref{p:guide}.
\end{proof}

The first half is O'Shaughnessy's thesis that trying is present in all action. That it often sounds odd to say that someone tried to do what she did with ease is a matter of what the remark suggests, not of what it says; it is true, and misleading, to say that I tried to open the unlocked door. The second half explains the structure of the intentions that ACT-11-a describes: the athlete intends to try to break the record, and her conduct is guided by breaking it, though she does not intend to break it.

\section{Objection and reply}

\begin{objection}
If guidance does not involve intention, trying collapses into wishing accompanied by activity.
\end{objection}

\begin{reply}
A wish need not guide anything; guidance requires that the agent's conduct be regulated by the aim, adjusted when it falls short, and abandoned when the aim is achieved or given up. This is what distinguishes the escapee who works at the bars from the prisoner who wishes he were free and paces his cell. The escapee intends to try; his trying is guided by escaping; neither the intention nor the guidance requires him to believe he will succeed. \citet{Grice1971} took intending to $A$ to involve believing that one will $A$; the account keeps that constraint for intention and locates the weaker, belief-independent attitude in the guidance of trying.
\end{reply}

\section{Conclusion}

To try to $A$ is to do something believed to bear on $A$-ing, under the guidance of $A$-ing. Guidance is weaker than intention, as it must be if the results of ACT-11-a are consistent. The account allows trying by mental acts alone, shows that all intentional action involves trying, and explains how an agent can intend to try what she believes she will not achieve.

\begin{thebibliography}{9}\small
\bibitem[Grice(1971)]{Grice1971} Grice, H.~P. (1971). Intention and uncertainty. \emph{Proceedings of the British Academy}, 57, 263--279.
\bibitem[Hornsby(1980)]{Hornsby1980} Hornsby, J. (1980). \emph{Actions}. Routledge \& Kegan Paul.
\bibitem[O'Shaughnessy(1973)]{OShaughnessy1973} O'Shaughnessy, B. (1973). Trying (as the mental \enquote{pineal gland}). \emph{The Journal of Philosophy}, 70(13), 365--386.
\end{thebibliography}

\end{document}
